% ============================================
% DOC-11
% Gerado de molambudos.md — build_miolo.py
% ============================================
\section*{\textbf{DOC-11} --- Diário de Oliveira I}
\markboth{DOC-11 --- Diário de Oliveira I}{}

\noindent \textit{O primeiro sintoma}

\noindent \textbf{Dr. Heitor Oliveira --- Caderno Pessoal}
\textbf{15 de abril de 1979}

\vspace{0.5em}\noindent\rule{\textwidth}{0.5pt}\vspace{0.5em}

\noindent \textit{Nota: Este caderno foi encontrado no armário do Dr. Oliveira após seu desaparecimento. As primeiras páginas são escritas em letra firme, de médico. As últimas são quase ilegíveis. A transcrição a seguir é literal.}

\vspace{0.5em}\noindent\rule{\textwidth}{0.5pt}\vspace{0.5em}

\noindent Dia 1.

\noindent A direção do hospital me deu uma semana de licença. Disse que eu \textquotedblleft{}precisava descansar\textquotedblright{}. O que eles querem dizer é que precisam de tempo para decidir se eu estou louco.

\noindent Não estou louco.

\noindent Mas estou com medo.

\noindent O medo é profissional: comecei a ter sintomas que diagnostico em meus pacientes. Insônia. Hipervigilância. Pensamentos intrusivos. O rosto do paciente 1.260 aparece na minha mente quando fecho os olhos. O olho amarelo.

\noindent Especificamente o olho.

\noindent Ele brilha.

\noindent Não é metáfora. Eu vi. Na cela, naquela noite, o olho do paciente brilhava no escuro como o olho de um gato. Não há explicação fisiológica. Não há doença que cause isso.

\noindent Tentei raciocinar: acúmulo de pigmento, reflexo anômalo da retina, uma variante rara de tapetum lucidum em humanos. Não. Não há registro. Não há literatura. Não há nada.

\noindent O olho brilha.

\noindent E agora, quando fecho os meus olhos, vejo ele.

\vspace{0.5em}\noindent\rule{\textwidth}{0.5pt}\vspace{0.5em}

\noindent Dia 2.

\noindent Acordei às 3:14 da manhã. O mesmo horário que o paciente 1.260 diz acordar todas as noites. Levantei, fui ao banheiro, lavei o rosto.

\noindent Quando olhei no espelho, meu olho direito parecia mais claro.

\noindent Fiquei 5 minutos olhando. Piscando. Abrindo e fechando. A luz do banheiro é branca, forte. Não há sombra. Não há ilusão.

\noindent A íris do olho direito está numa tonalidade mais clara que a do esquerdo.

\noindent Isso é impossível. A cor dos olhos não muda na idade adulta. Sinto que estou perdendo o juízo.

\noindent Mas não estou.

\noindent Estou registrando.

\noindent Estou documentando.

\noindent Porque, se eu estiver ficando louco, quero que alguém leia isto e saiba \textit{como} aconteceu.

\vspace{0.5em}\noindent\rule{\textwidth}{0.5pt}\vspace{0.5em}

\noindent Dia 3.

\noindent O cheiro.

\noindent Acordei e o quarto estava cheio de um cheiro adocicado. Como carne doce. Como flor murcha. O mesmo cheiro que senti na cela do paciente 1.260.

\noindent Levantei, abri a janela. O ar externo estava limpo, frio, normal. Mas dentro do quarto, o cheiro persistia.

\noindent Estava no travesseiro. No lençol. Na minha pele.

\noindent Tomei banho. Esfreguei a pele com sabão. O cheiro saiu --- mas voltou minutos depois. Como se estivesse saindo de dentro de mim.

\noindent Lavei a roupa de cama. O cheiro saiu.

\noindent Quando deitei, o cheiro voltou.

\noindent Não é o ambiente. Sou eu.

\noindent Estou exalando o cheiro.

\noindent Fiz um experimento: pedi para minha esposa cheirar meu braço. Ela disse: \textquotedblleft{}Não sinto nada.\textquotedblright{} Então é subjetivo? Olfativo? Alucinatório?

\noindent Pedi para ela cheirar de novo, com atenção. Ela disse: \textquotedblleft{}Tem um cheiro... estranho? Doce? Agora que você falou, talvez.\textquotedblright{}

\noindent Ela sentiu.

\noindent Não é alucinação.

\noindent É contágio.

\vspace{0.5em}\noindent\rule{\textwidth}{0.5pt}\vspace{0.5em}

\noindent Dia 4.

\noindent Sonhei com a vala.

\noindent Estava deitado na pilha de corpos. Os olhos deles estavam abertos. A criatura estava em cima de mim. O rosto perto do meu. Ela disse:

\noindent \textit{\textquotedblleft{}Você leu. Agora você é meu.\textquotedblright{}}

\noindent Acordei gritando.

\noindent Minha esposa dorme no sofá agora. Disse que eu falo dormindo. Perguntei o que eu falo. Ela não quis repetir. Insisti.

\noindent Ela disse: \textquotedblleft{}Você diz: 'O ciclo recomeça. Paciente 1.261.'\textquotedblright{}

\noindent Paciente 1.261.

\noindent O número depois do 1.260.

\noindent \textbf{Eu.}

\vspace{0.5em}\noindent\rule{\textwidth}{0.5pt}\vspace{0.5em}

\noindent Dia 5.

\noindent Fui ao hospital. Entrei na ala dos silenciados. Fiquei na porta da cela 1.260. O paciente não está mais lá --- transferiram para a enfermaria geral depois que as sessões pararam.

\noindent Mas a cela está vazia.

\noindent Entrei.

\noindent O cheiro adocicado é forte aqui. Mais forte que em qualquer outro lugar.

\noindent No chão, perto da cama, havia um papel. Peguei. Estava escrito, em caligrafia minúscula:

\noindent \textit{\textquotedblleft{}O médico vai voltar. Ele já sonhou com a vala. Ele já sentiu o cheiro. Ele já é meu.\textquotedblright{}}

\noindent A data no papel: 1919.

\noindent Escrito há 60 anos.
Prevendo que eu estaria aqui.
Prevendo tudo.

\noindent Saí correndo da cela. Passei pelo corredor. Os pacientes olharam. Alguns riram. Outros gritaram.

\noindent Um deles, um velho na cela do fundo, disse quando passei:

\noindent \textit{\textquotedblleft{}Bem-vindo, doutor. A criatura estava esperando o senhor.\textquotedblright{}}

\noindent Parei.

\noindent \textit{\textquotedblleft{}O que você disse?\textquotedblright{}}

\noindent \textit{\textquotedblleft{}Ela disse. Disse que o senhor ia vir. Disse que o senhor ia sentir o cheiro. Disse que o senhor ia sonhar.\textquotedblright{}}

\noindent \textit{\textquotedblleft{}Quem disse?\textquotedblright{}}

\noindent \textit{\textquotedblleft{}A Rasga Mortalha.\textquotedblright{}}

\noindent Ele apontou para a janela.

\noindent No batente, uma coruja branca me olhava.

\noindent Quando pisquei, ela tinha ido embora.

\noindent Mas o bilhete de 1919 está no meu bolso.

\noindent E a letra --- a letra está mudando.

\noindent A minha letra.

\noindent Está ficando igual à dele.

\vspace{0.5em}
\noindent\textit{\textbf{↪ Rotas:} \rota{DOC-12} \rotasep \rota{DOC-06} \rotasep \rota{LUC-03}}